\documentclass[UTF8,12pt]{ctexart}
\usepackage[a4paper,margin=24mm]{geometry}
\usepackage{amsmath,amssymb,bm,graphicx,adjustbox}
\newcommand{\fitmath}[1]{\begin{adjustbox}{max width=\linewidth}$\displaystyle #1$\end{adjustbox}}
\setlength{\parindent}{0pt}
\setlength{\parskip}{.7em}
\sloppy
\allowdisplaybreaks
\begin{document}
\section*{2010 年考研数学三 · 真题}
\subsection*{原 PDF 第 46 页}

\subsection*{2010年全国硕士研究生招生考试试题}

\subsection*{一、选择题（本题共8小题，每小题4分，共32分。在每小题给出的四个选项中，只有一项符合题目要求，把所选项前的字母填在题后的括号内。）}

（1）若 \(\displaystyle \lim_{x\to0}\left[\dfrac{\displaystyle 1}{\displaystyle x}-\left(\dfrac{\displaystyle 1}{\displaystyle x}-a\right)e^x\right]=1\)，则 \(\displaystyle a\) 等于（　　）
（A）\(\displaystyle 0\)。 （B）\(\displaystyle 1\)。 （C）\(\displaystyle 2\)。 （D）\(\displaystyle 3\)。


（2）设 \(\displaystyle y_1,\allowbreak y_2\) 是一阶线性非齐次微分方程 \(\displaystyle y^{\prime}+p(x)y=q(x)\) 的两个特解，若常数 \(\displaystyle \lambda,\allowbreak \mu\) 使 \(\displaystyle \lambda y_1+\mu y_2\) 是该方程的解，\(\displaystyle \lambda y_1-\mu y_2\) 是该方程对应的齐次方程的解，则（　　）
（A）\(\displaystyle \lambda=\dfrac{\displaystyle 1}{\displaystyle 2},\allowbreak \mu=\dfrac{\displaystyle 1}{\displaystyle 2}\)。 （B）\(\displaystyle \lambda=-\dfrac{\displaystyle 1}{\displaystyle 2},\allowbreak \mu=-\dfrac{\displaystyle 1}{\displaystyle 2}\)。 （C）\(\displaystyle \lambda=\dfrac{\displaystyle 2}{\displaystyle 3},\allowbreak \mu=\dfrac{\displaystyle 1}{\displaystyle 3}\)。 （D）\(\displaystyle \lambda=\dfrac{\displaystyle 2}{\displaystyle 3},\allowbreak \mu=\dfrac{\displaystyle 2}{\displaystyle 3}\)。


（3）设函数 \(\displaystyle f(x),\allowbreak g(x)\) 具有二阶导数，且 \(\displaystyle g^{\prime\prime}(x)<0\)。若 \(\displaystyle g(x_0)=a\) 是 \(\displaystyle g(x)\) 的极值，则 \(\displaystyle f(g(x))\) 在 \(\displaystyle x_0\) 处取极大值的一个充分条件是（　　）
（A）\(\displaystyle f^{\prime}(a)<0\)。 （B）\(\displaystyle f^{\prime}(a)>0\)。 （C）\(\displaystyle f^{\prime\prime}(a)<0\)。 （D）\(\displaystyle f^{\prime\prime}(a)>0\)。


（4）设 \(\displaystyle f(x)=\ln^{10}x\)，\(\displaystyle g(x)=x\)，\(\displaystyle h(x)=e^{x/10}\)，则当 \(\displaystyle x\) 充分大时有（　　）
（A）\(\displaystyle g(x)<h(x)<f(x)\)。 （B）\(\displaystyle h(x)<g(x)<f(x)\)。 （C）\(\displaystyle f(x)<g(x)<h(x)\)。 （D）\(\displaystyle g(x)<f(x)<h(x)\)。


（5）设向量组I：\(\displaystyle \boldsymbol\alpha_1,\allowbreak \boldsymbol\alpha_2,\allowbreak \ldots,\allowbreak \boldsymbol\alpha_r\) 可由向量组II：\(\displaystyle \boldsymbol\beta_1,\allowbreak \boldsymbol\beta_2,\allowbreak \ldots,\allowbreak \boldsymbol\beta_s\) 线性表示。下列命题正确的是（　　）
（A）若向量组I线性无关，则 \(\displaystyle r\le s\)。 （B）若向量组I线性相关，则 \(\displaystyle r>s\)。
（C）若向量组II线性无关，则 \(\displaystyle r\le s\)。 （D）若向量组II线性相关，则 \(\displaystyle r>s\)。


（6）设 \(\displaystyle A\) 为4阶实对称矩阵，且 \(\displaystyle A^2+A=O\)。若 \(\displaystyle A\) 的秩为3，则 \(\displaystyle A\) 相似于（　　）
（A）\(\displaystyle \operatorname{diag}(1,\allowbreak 1,\allowbreak 1,\allowbreak 0)\)。 （B）\(\displaystyle \operatorname{diag}(1,\allowbreak 1,\allowbreak -1,\allowbreak 0)\)。
（C）\(\displaystyle \operatorname{diag}(1,\allowbreak -1,\allowbreak -1,\allowbreak 0)\)。 （D）\(\displaystyle \operatorname{diag}(-1,\allowbreak -1,\allowbreak -1,\allowbreak 0)\)。


\clearpage

\subsection*{原 PDF 第 47 页}

（7）设随机变量 \(\displaystyle X\) 的分布函数


\[\fitmath{\displaystyle F(x)=\begin{cases}0,&x<0,\\\dfrac{\displaystyle 1}{\displaystyle 2},&0\le x<1,\\1-e^{-x},&x\ge1,\end{cases}}\]


则 \(\displaystyle P\{X=1\}=\)（　　）
（A）\(\displaystyle 0\)。 （B）\(\displaystyle \dfrac{\displaystyle 1}{\displaystyle 2}\)。 （C）\(\displaystyle \dfrac{\displaystyle 1}{\displaystyle 2}-e^{-1}\)。 （D）\(\displaystyle 1-e^{-1}\)。


（8）设 \(\displaystyle f_1(x)\) 为标准正态分布的概率密度，\(\displaystyle f_2(x)\) 为 \(\displaystyle [-1,\allowbreak 3]\) 上均匀分布的概率密度，若


\[\fitmath{\displaystyle f(x)=\begin{cases}af_1(x),&x\le0,\\bf_2(x),&x>0\end{cases}\quad(a>0,b>0)}\]


为概率密度，则 \(\displaystyle a,\allowbreak b\) 应满足（　　）
（A）\(\displaystyle 2a+3b=4\)。 （B）\(\displaystyle 3a+2b=4\)。 （C）\(\displaystyle a+b=1\)。 （D）\(\displaystyle a+b=2\)。


\subsection*{二、填空题（本题共6小题，每小题4分，共24分，把答案填在题中横线上。）}

（9）设可导函数 \(\displaystyle y=y(x)\) 由方程 \(\displaystyle \int_0^{x+y}e^{-t^2}\,dt=\displaystyle \int_0^x x\sin t^2\,dt\) 确定，则 \(\displaystyle \left.\dfrac{\displaystyle dy}{\displaystyle dx}\right|_{x=0}=\)\_\_\_\_\_\_。


（10）设位于曲线 \(\displaystyle y=\dfrac{\displaystyle 1}{\displaystyle x\sqrt{1+\ln^2x}}\; (e\le x<+\infty)\) 下方、\(\displaystyle x\) 轴上方的无界区域为 \(\displaystyle G\)，则 \(\displaystyle G\) 绕 \(\displaystyle x\) 轴旋转一周所得空间区域的体积为\_\_\_\_\_\_。


（11）设某商品的收益函数为 \(\displaystyle R(p)\)，收益弹性为 \(\displaystyle 1+p^3\)，其中 \(\displaystyle p\) 为价格，且 \(\displaystyle R(1)=1\)，则 \(\displaystyle R(p)=\)\_\_\_\_\_\_。


（12）若曲线 \(\displaystyle y=x^3+ax^2+bx+1\) 有拐点 \(\displaystyle (-1,\allowbreak 0)\)，则 \(\displaystyle b=\)\_\_\_\_\_\_。


（13）设 \(\displaystyle A,\allowbreak B\) 为3阶矩阵，且 \(\displaystyle |A|=3,\allowbreak \ |B|=2,\allowbreak \ |A^{-1}+B|=2\)，则 \(\displaystyle |A+B^{-1}|=\)\_\_\_\_\_\_。


（14）设 \(\displaystyle X_1,\allowbreak X_2,\allowbreak \ldots,\allowbreak X_n\) 是来自总体 \(\displaystyle N(\mu,\allowbreak \sigma^2)\; (\sigma>0)\) 的简单随机样本。记统计量 \(\displaystyle T=\dfrac{\displaystyle 1}{\displaystyle n}\displaystyle \sum_{i=1}^{n}X_i^2\)，则 \(\displaystyle E(T)=\)\_\_\_\_\_\_。


\subsection*{三、解答题（本题共9小题，共94分，解答应写出文字说明、证明过程或演算步骤。）}

（15）（本题满分10分）求极限 \(\displaystyle \lim_{x\to+\infty}(x^{1/x}-1)^{1/\ln x}\)。


\clearpage

\subsection*{原 PDF 第 48 页}

（16）（本题满分10分）计算二重积分 \(\displaystyle \iint_D(x+y)^3\,dx\,dy\)，其中 \(\displaystyle D\) 由曲线 \(\displaystyle x=\sqrt{1+y^2}\) 与直线 \(\displaystyle x+\sqrt2y=0\) 及 \(\displaystyle x-\sqrt2y=0\) 围成。


（17）（本题满分10分）求函数 \(\displaystyle u=xy+2yz\) 在约束条件 \(\displaystyle x^2+y^2+z^2=10\) 下的最大值和最小值。


（18）（本题满分10分）
（I）比较 \(\displaystyle \int_0^1|\ln t|[\ln(1+t)]^n\,dt\) 与 \(\displaystyle \int_0^1t^n|\ln t|\,dt\; (n=1,\allowbreak 2,\allowbreak \ldots)\) 的大小，说明理由；
（II）记 \(\displaystyle u_n=\displaystyle \int_0^1|\ln t|[\ln(1+t)]^n\,dt\; (n=1,\allowbreak 2,\allowbreak \ldots)\)，求极限 \(\displaystyle \lim_{n\to\infty}u_n\)。


（19）（本题满分10分）设函数 \(\displaystyle f(x)\) 在 \(\displaystyle [0,\allowbreak 3]\) 上连续，在 \(\displaystyle (0,\allowbreak 3)\) 内存在二阶导数，且 \(\displaystyle 2f(0)=\displaystyle \int_0^2f(x)\,dx=f(2)+f(3)\)。
（I）证明存在 \(\displaystyle \eta\in(0,\allowbreak 2)\)，使 \(\displaystyle f(\eta)=f(0)\)；
（II）证明存在 \(\displaystyle \xi\in(0,\allowbreak 3)\)，使 \(\displaystyle f^{\prime\prime}(\xi)=0\)。


\clearpage

\subsection*{原 PDF 第 49 页}

（20）（本题满分11分）设


\[\fitmath{\displaystyle A=\begin{pmatrix}\lambda&1&1\\0&\lambda-1&0\\1&1&\lambda\end{pmatrix},\quad\boldsymbol b=\begin{pmatrix}a\\1\\1\end{pmatrix}.}\]


已知线性方程组 \(\displaystyle A\boldsymbol x=\boldsymbol b\) 存在两个不同的解。
（I）求 \(\displaystyle \lambda,\allowbreak a\)；
（II）求方程组 \(\displaystyle A\boldsymbol x=\boldsymbol b\) 的通解。


（21）（本题满分11分）设


\[\fitmath{\displaystyle A=\begin{pmatrix}0&-1&4\\-1&3&a\\4&a&0\end{pmatrix},}\]


正交矩阵 \(\displaystyle Q\) 使 \(\displaystyle Q^{\mathrm T}AQ\) 为对角矩阵，若 \(\displaystyle Q\) 的第1列为 \(\displaystyle \dfrac{\displaystyle 1}{\displaystyle \sqrt6}(1,\allowbreak 2,\allowbreak 1)^{\mathrm T}\)，求 \(\displaystyle a,\allowbreak Q\)。


（22）（本题满分11分）设二维随机变量 \(\displaystyle (X,\allowbreak Y)\) 的概率密度为 \(\displaystyle f(x,\allowbreak y)=Ae^{-2x^2+2xy-y^2},\allowbreak \ -\infty<x<+\infty,\allowbreak \ -\infty<y<+\infty\)，求常数 \(\displaystyle A\) 及条件概率密度 \(\displaystyle f_{Y\mid X}(y\mid x)\)。


（23）（本题满分11分）箱中装有6个球，其中红、白、黑球的个数分别为1，2，3个。现从箱中随机地取出2个球，记 \(\displaystyle X\) 为取出的红球个数，\(\displaystyle Y\) 为取出的白球个数。
（I）求随机变量 \(\displaystyle (X,\allowbreak Y)\) 的概率分布；
（II）求 \(\displaystyle \operatorname{Cov}(X,\allowbreak Y)\)。


\clearpage
\end{document}
